
au degré @@M0@@.

Le module des différentielles @@M1@@ est isomorphe à @@M2@@ car on peut définir une @@M3@@-dérivation @@M4@@ :
@@M5@@
de façon à satisfaire la propriété universelle qui définit @@M6@@ en posant @@M7@@ ; de plus, @@M8@@ est compatible aux graduations si on attribue le degré @@M9@@ au facteur @@M10@@.

Soit @@M11@@ le @@M12@@-fibré vectoriel défini par @@M13@@. Alors @@M14@@ est la restriction à @@M15@@ de @@M16@@. Or @@M17@@ est le @@M18@@-schéma affine défini par @@M19@@, donc @@M20@@ est défini par @@M21@@ (car @@M22@@ est localement libre de type fini), donc @@M23@@ est égal à @@M24@@, et par suite, @@M25@@. Pour voir sa graduation, on regarde localement, et on voit aussitôt que @@M26@@ doit être considéré comme étant de degré @@M27@@.

En prenant les termes de degré @@M28@@ dans la suite exacte (5.2.1), on trouve la suite exacte annoncée.

\textbf{Remarque 5.2.1.} En utilisant les résultats généraux sur l'étude différentielle des foncteurs de Hilbert ([2], 5), on obtient directement un isomorphisme @@M29@@ d'après la caractérisation fonctorielle donnée de @@M30@@ pour tout @@M31@@-morphisme @@M32@@.

\textbf{5.3.} D'après 5.2, on a :
@@M33@@
Par conséquent :
@@M34@@

